% !PRINTED 321
\weberpage{321}
%<<CONT
versely show without difficulty that the four quantities
$X$ contained in the form (8), when the $v_{1,0}, v_{0,1}, v_{1,1}$ are the roots of any
cubic equation, satisfy a biquadratic equation with
rational coefficients\textsuperscript{1)}.

\noindent\rule{0.35\linewidth}{0.4pt}

\footnotesize
\noindent\textsuperscript{1)} This form of the solution of the biquadratic equation is the
analogue of the \gesperrt{Cayley} form of the \gesperrt{Cardan formula} (Bd. I,
\S.\,36). I do not know whether this form of the solution of biquadratic
equations is otherwise already known in elementary algebra. I
first found it in an English collection of problems ``Mathematical
Tripos, Part I, June 1894'', where the cubic resolvent
of the $v$ is also formed.
\normalsize
